\subsection{A Calabi--Eckmann family without LCK metrics}
\label{sec:calabi-eckmann-endpoints}
\index{Calabi--Eckmann manifold!determinant endpoint}

A second construction separates the endpoint criterion from LCK geometry. Let $\tau=A+\mathrm iB_0$, where $A<0$ and $B_0>0$, and let $\C$ act holomorphically on $(\C^2\setminus\{0\})^2$ by
\[
 t\cdot(z,w)=(\mathrm e^t z,\mathrm e^{\tau t}w).
\]
Write $X_\tau$ for the quotient and put $r_1=\log|z|^2$, $r_2=\log|w|^2$. On the projective base let $\omega_1,\omega_2$ be the normalized pulled Fubini--Study forms.

\begin{proposition}\label{prop:calabi-eckmann-endpoint}
The quotient is a compact complex threefold diffeomorphic to $S^3\times S^3$. Projectivization gives a principal elliptic fibration over $\mathbb P^1\times\mathbb P^1$. The form $\eta=\tau\partial r_1-\partial r_2$ descends, and for
\[
 g=\pi^*(\omega_1+\omega_2)+s\mathrm i\eta\wedge\bar\eta,
 \qquad s>0,
\]
one has
\begin{equation}\label{eq:calabi-eckmann-lower-mass}
 v_-(g)=48\pi^2sB_0.
\end{equation}
Every Hermitian $h\geq c g$ satisfies
\[
 \sup_{T\in\mathscr G_h}\|k_T\|_{L^{3/2}(\mathrm dV)}
 \leq c^{-1}(48\pi^2sB_0)^{-1/3}.
\]
No balanced, pluriclosed or LCK metric exists on $X_\tau$.
\end{proposition}
\begin{proof}
Every orbit has a unique point with $|z|=|w|=1$: the needed parameter solves
\[
 \operatorname{Re}t=-r_1/2,\qquad
 \operatorname{Im}t=(r_2-Ar_1)/(2B_0).
\]
These smooth formulas give a product trivialization by $\C\times(S^3\times S^3)$, proving freeness and properness. A branch of the logarithm sets one nonzero $z$-coordinate to one, with local holomorphic coordinates $z_{\rm other}/z_{\rm chosen}$ and $w_jz_{\rm chosen}^{-\tau}$. Thus the quotient carries the asserted compact complex structure. Logarithm branches account for the residual fiber periods.

The invariant form $\eta$ annihilates the holomorphic action generator, so it descends. On an elliptic fiber it is $\mathrm d\chi$, where $\chi=\tau\log v_1-\log v_2$ has lattice generated by $2\pi\mathrm i\tau$ and $-2\pi\mathrm i$. Hence $A_E=8\pi^2B_0$. Direct differentiation gives
\[
 \begin{aligned}
 F&=2\pi\mathrm i(\tau\omega_1-\omega_2),&
 \alpha&=-2\pi B_0\omega_1,\\
 \beta&=2\pi(A\omega_1-\omega_2),&
 Q&=-8\pi^2A\,\omega_1\wedge\omega_2>0.
 \end{aligned}
\]
Since $\int_B(\omega_1+\omega_2)^2=2$, \cref{thm:elliptic-curvature-endpoint} gives the mass, endpoint and pluriclosed obstruction.

Moreover $\mathrm d\operatorname{Re}\eta=-2\pi B_0\pi^*\omega_1$. This nonzero semipositive exact form excludes every balanced metric by the same Stokes pairing as above, hence also every K\"ahler metric. The quotient is simply connected. A hypothetical LCK metric therefore has exact Lee form $\mathrm df$; multiplying it by $\mathrm e^{-f}$ gives a K\"ahler metric, a contradiction. In particular no Vaisman metric exists.
\end{proof}

The curvature criterion controls bounded potentials and the full Gauduchon current family in these threefold geometries. Changing-sign curvature, primitive elliptic bundles and higher-dimensional analogues require further estimates.
